% 05_results_B.tex — Results B: attribution / XAI. Drafted by item P2-E8-6 (P2 Sprint 7).
% Voice: ../STYLE.md (Andreas Maier). Numbers: tools/xai_study/compare/out/leaderboard.md
% + faithful_demo.md + the §R per-method records under tools/xai_study/phaseB_attribution/out/.
% Honesty contract: ../plan.md — faithfulness only; no semantics/meaning recovered here.

\section{Results B --- attributing a VCS output to its causes}\label{sec:results-b}

Attribution is the field's everyday tool: hand it a system and an output, and it
hands back a map of which inputs mattered. We turn the toolkit on the VCS and ask
the one question a real network never lets us ask --- is the map \emph{true}? Every
attribution we compute is scored against the same \S1 intervention oracle that
defines ground truth (Sec.~\ref{sec:methods}), so for each chosen output $y$ ---
a pixel, a score, a game event --- we know the exact set of causes the explanation
should have named. We ran twelve attribution methods across six core games (Pong,
Breakout, Space~Invaders, Seaquest, Ms.~Pac-Man, Q*bert) and an oracle-as-method
positive control. Figure~\ref{fig:fig4} (Phase-B half) places them on the shared
faithfulness axis; Table~\ref{tab:phaseB} reports the per-method scores.

The headline is a split, not a ranking. The VCS computes two kinds of output, and
the gradient is alive for one and dead for the other. A \emph{content} output ---
the value of a register, a colour, a graphics bit --- is a differentiable function
of the state, so its gradient flows. A \emph{position} output --- where a sprite
is drawn, set by strobe timing on a discrete index --- has \emph{provably zero}
naive gradient (Paper~1, and Sec.~\ref{sec:methods}). The entire gradient family
inherits this boundary, and it cuts the leaderboard cleanly in two.

\begin{figure}[t]
\centering
\includegraphics[width=\textwidth]{fig4_attribution_vs_mechanistic.pdf}
\caption{\textbf{Attribution (Phase~B) and mechanistic interpretability
(Phase~C) scored against the same \S1 intervention oracle.} The Phase-B half (a)
splits each method's faithfulness by the output regime it targets: on a
\emph{content} output (a register/colour/graphics-bit value, a differentiable
function of state) the gradient family lands on the true cause, but on a
\emph{position}/index output the entire gradient family --- vanilla saliency,
SmoothGrad, guided backprop, Grad$\times$Input/DeepLIFT, Integrated Gradients,
Expected Gradients, and the gradient-leaning surrogates LIME, KernelSHAP, RISE ---
\emph{collapses to $0.00$}, while the intervention methods (occlusion, extremal
perturbation, on-distribution counterfactual) hold up because each probe is a real
re-run of the bit-exact program. \emph{Note that} the popular gradient methods,
which fail on position, carry the higher human-plausibility proxy --- the danger
zone of Sec.~\ref{sec:results_compare}, in numbers. The Phase-C half (b) is
discussed in Sec.~\ref{sec:results-c}. Vector figure; values read from the
committed cross-tradition leaderboard.}\label{fig:fig4}
\end{figure}

\subsection*{The gradient family: faithful on content, blind on position}

On the content path the gradient is faithful and, frankly, unremarkable. Vanilla
saliency~\cite{simonyan2014deep} lands its mass on the true causal byte: on
Q*bert its attribution of a content read correlates with the oracle at Pearson
$0.999$ with precision@3 of $1.0$, on Ms.~Pac-Man at $0.811$ / $1.0$, on Breakout
at $0.750$ / $1.0$. Integrated Gradients~\cite{sundararajan2017axiomatic},
Grad$\times$Input/DeepLIFT~\cite{shrikumar2017learning},
SmoothGrad~\cite{smilkov2017smoothgrad}, and guided
backpropagation~\cite{springenberg2015striving} all reproduce this, because the
output we explain is a one-hot read whose Jacobian $\partial y/\partial u$ is a
constant pointing at the read index. Smoothing denoises the picture but does not
move the score ($\Delta\,\mathrm{corr}\approx 0$); guided backprop degenerates to
vanilla saliency, since the soft VCS has no deep ReLU stack for its
negative-gradient suppression to act on. The methods agree because, on this path,
there is nothing to disagree about.

The position output is where the family fails, and it fails by vanishing. For the
discrete sprite-position output the naive gradient is identically zero on every
game, so saliency, Grad$\times$Input, SmoothGrad, guided backprop, and Integrated
Gradients all collapse to a faithfulness of $0.0$ in the position regime --- the
``Fail / zero on index outputs'' prediction of Sec.~\ref{sec:methods}, now a
measured number. Averaged over all output regimes the gradient family scores only
$0.294$ faithfulness against the oracle (Fig.~\ref{fig:fig4}); on the position
regime alone, $0.068$. A method that returns the empty map for half the outputs a
game produces is not a partial success but a structural blind spot, and the spot
is exactly where the discrete game logic lives.

Two finer results sharpen the warning rather than soften it. First, the choice of
baseline --- the perennial worry about Integrated Gradients --- turns out to be a
\emph{magnitude} effect here, not a faithfulness one: across the baseline sweep the
$(x-x_0)$ factor moves $\max|\mathrm{IG}|$, but the ranking and the correlation
with the oracle are invariant, because the constant one-hot Jacobian keeps all the
mass on the same byte. The only baseline that breaks IG is the degenerate one, a
baseline byte equal to the content byte, which kills the attribution outright.
Expected Gradients~\cite{erion2021improving}, averaging IG over a distribution of
real recorded VCS states, removes that degeneracy and is provably stable across
reference draws (mean pairwise cosine $\approx 1$) --- yet on the content path,
where $\partial y/\partial u$ is constant, it scores the \emph{same} faithfulness
as a well-chosen single-baseline IG at $M\times$ the compute, and it still
vanishes on position. On distribution, the more careful gradient method buys
stability, not truth; its all-regime faithfulness of $0.034$ is the lowest in the
study.

Second, and most telling, guided backprop and vanilla saliency fail the
program-randomization sanity check of Adebayo et al.~\cite{adebayo2018sanity} ---
on the system itself, in numbers. We scramble the ROM so the program boots to a
completely different state ($\sim$95--127 of 128 RAM bytes change), and the
content-path saliency map stays \emph{byte-identical} (cosine $\approx 1.0$ on
all six games). The Jacobian $\partial y/\partial u = e_{\mathrm{idx}}$ is a
constant one-hot, independent of what the program computes, so the map points at
the read index no matter which program produced the output --- a saliency that is
\emph{provably model-invariant}. A faithful explanation must change when the model
changes; this one cannot. By contrast the intervention oracle, which re-runs the
real program for every candidate cause, distinguishes the two programs at once.
This is the danger zone made concrete: a gradient map that lands on the right
cause and is still blind to the model that produced it.

\subsection*{The intervention family: faithful on both regimes}

Replace the gradient with an actual intervention and the position blind spot
disappears. Occlusion~\cite{zeiler2014visualizing} sets a candidate cause and
re-runs the bit-exact program, so it is the coarse twin of the oracle itself; on
the Pong position output, where every gradient method scores $0.0$, occlusion
recovers the true causal bytes at correlation $0.837$. Its all-regime faithfulness
of $0.482$ is the highest in Phase~B. The black-box perturb-and-re-run methods
follow the same logic and the same result: RISE~\cite{petsiuk2018rise} ($0.375$),
LIME~\cite{ribeiro2016why} ($0.382$), and KernelSHAP / Shapley
sampling~\cite{lundberg2017unified,strumbelj2014explaining} ($0.446$) each draw
masked coalitions, re-run the real ROM, and fit or average over the genuine
responses, so each works on position outputs where the gradient is dead. They
converge as the sample budget $N$ grows --- their variance falls as $1/N$, the
expected approximation behaviour --- and they report their own completeness
($\Sigma\phi = y(\text{intact}) - y(\text{all-masked})$ for KernelSHAP, the local
surrogate $R^2$ for LIME).

The two perturbation methods built to find a \emph{minimal} mask close the loop on
Atrey et al.'s~\cite{atrey2020exploratory} objection that perturbations stray
off-manifold. Extremal perturbation~\cite{fong2017interpretable,fong2019understanding}
optimises an area-bounded high-impact mask in which every probe is a real
$do(\text{occlude})$ re-run, and the on-distribution
counterfactual~\cite{olson2021counterfactual} substitutes candidate cells toward a
\emph{genuine} alternative VCS state --- the RAM of another frame of the same ROM
--- before re-running. Both stay on the data manifold by construction, both score
a true-causal-set mask IoU and a minimality ratio against the oracle, and both
work on position outputs (faithfulness $0.346$ and $0.350$) precisely because they
re-run the program rather than differentiate it. The intervention family as a
whole averages $0.681$ faithfulness against the oracle, more than double the
gradient family's $0.294$ (Fig.~\ref{fig:fig4}). The lesson is mechanistic, not
incidental: a perturbation is faithful exactly when it is a valid intervention.

\subsection*{Methods that do not even apply}

A part of the popular toolkit returns no answer at all, because it has nothing on
the VCS to attach to. Grad-CAM and Grad-CAM++~\cite{selvaraju2017gradcam,%
chattopadhyay2018gradcampp} need a final convolutional layer to weight; attention
rollout~\cite{abnar2020quantifying} needs attention matrices;
VIPER~\cite{bastani2018verifiable} needs a learned policy to distil. The VCS is a
6502 CPU with a TIA and a RIOT executing a fixed ROM --- no conv channels, no
attention heads, no policy network --- so across all six games the count of each
required substrate is exactly zero, against $\geq 16$ genuine candidate causes per
game for the applicable methods. We record this as a measured finding, $\mathtt{applies
= false}$, not a poor score: the absence is structural, and it is the same in every
game. Re-targeting Grad-CAM's pooling onto raw register state would not be
Grad-CAM, only the content-path gradient under a new name. That a slice of the
field's headline methods is silent on a real, deployed, deterministic artifact is
itself part of the audit --- the measured narrowness of architecture-specific
attribution.

\subsection*{Plausible is not faithful}

Across the twelve methods the leaderboard separates cleanly along the axis that
matters and crosses the axis the field actually optimizes. The single starkest
pair makes it concrete: on the position regime the causal method activation
patching (Phase~C) reaches faithfulness $1.000$, the oracle ceiling, on all six
games, while vanilla saliency --- the field's default tool --- collapses to
$0.000$ on all six; a per-method gap of $1.000$, and bit-exact on every game
individually. Yet on the human-plausibility proxy of Sec.~\ref{sec:methods} the
\emph{popular} method carries the \emph{higher} score ($0.9$ versus $0.5$): the
map that is provably empty looks more convincing than the one that is exactly
right. That inversion --- low faithfulness paired with high plausibility --- is
the danger zone, and on the VCS it is not a worry but a measurement. The most
plausible Phase-B methods (the gradient family, plausibility $0.9$) include the
four with zero position-regime faithfulness; the most faithful (the intervention
family) wear the lower plausibility tag.

We claim no more than the faithfulness of these maps. A correct occlusion map
names the true causal bytes; it does not tell us what those bytes \emph{mean} ---
that a cell is the ball's $x$, that a routine is the bounce. Object-level
semantics enter only through the external T3 labels (AtariARI, OCAtari;
Sec.~\ref{sec:methods}), and even there we \emph{verify} a supplied label by
intervention rather than \emph{discover} it. Whether a faithful attribution
amounts to an understanding of the computation is the question we take up in
Sec.~\ref{sec:results_compare}; here we have established only that faithfulness
can be had --- and that the most popular way of seeking it does not deliver it.

\begin{table}[t]
\caption{\textbf{Phase-B attribution scored against the \S1 intervention oracle}
(six core games; oracle-as-method positive control $=1.0$ on every non-degenerate
column). \emph{Faith} is the all-regime faithfulness against ground truth and
\emph{Faith (pos)} the position/index regime alone, where the naive gradient is
provably zero; \emph{Plaus} is the documented method-tradition plausibility
\emph{proxy} (Sec.~\ref{sec:methods}), \emph{not} a human measurement.
\emph{Note that} the entire gradient family scores $0.000$ on position while the
intervention family does not, and that the popular gradient methods carry the
\emph{higher} plausibility tag --- the danger zone, measured. Values are read from
the committed cross-tradition leaderboard.}\label{tab:phaseB}
\begin{tabular}{@{}llccc@{}}
\toprule
Method & Tradition & Faith & Faith (pos) & Plaus (proxy) \\
\midrule
Occlusion~\cite{zeiler2014visualizing}            & intervention & 0.482 & 0.306 & 0.55 \\
KernelSHAP~\cite{lundberg2017unified}             & gradient/black-box & 0.446 & 0.279 & 0.90 \\
LIME~\cite{ribeiro2016why}                        & gradient/black-box & 0.382 & 0.167 & 0.90 \\
RISE~\cite{petsiuk2018rise}                        & gradient/black-box & 0.375 & 0.169 & 0.90 \\
Grad$\times$Input / DeepLIFT~\cite{shrikumar2017learning} & gradient & 0.367 & 0.000 & 0.90 \\
On-distribution counterfactual~\cite{olson2021counterfactual} & intervention & 0.350 & 0.139 & 0.55 \\
Extremal perturbation~\cite{fong2019understanding} & intervention & 0.346 & 0.202 & 0.55 \\
Integrated Gradients~\cite{sundararajan2017axiomatic} & gradient & 0.279 & 0.000 & 0.90 \\
Guided backprop~\cite{springenberg2015striving}   & gradient & 0.267 & 0.000 & 0.90 \\
SmoothGrad~\cite{smilkov2017smoothgrad}           & gradient & 0.267 & 0.000 & 0.90 \\
Vanilla saliency~\cite{simonyan2014deep}          & gradient & 0.267 & 0.000 & 0.90 \\
Expected Gradients~\cite{erion2021improving}      & gradient & 0.034 & 0.000 & 0.90 \\
\midrule
\multicolumn{2}{@{}l}{\emph{Family means} --- gradient/correlational} & 0.294 & 0.068 & --- \\
\multicolumn{2}{@{}l}{\hphantom{\emph{Family means} ---} intervention} & 0.681 & 0.412 & --- \\
Grad-CAM / attention / VIPER~\cite{selvaraju2017gradcam,abnar2020quantifying,bastani2018verifiable}
 & N/A & \multicolumn{3}{c}{does not apply (no conv / attention / policy)} \\
\botrule
\end{tabular}
\end{table}

% =====================================================================
% BIBTEX-NEEDED: full entries for the keys this section cites that are
% NOT yet in references.bib (SM to merge; do NOT edit references.bib here).
% Already-present keys reused as-is (NO entry below): sundararajan2017axiomatic,
% selvaraju2017gradcam, chattopadhyay2018gradcampp, atrey2020exploratory.
% =====================================================================
%
% @inproceedings{simonyan2014deep,
%   title={Deep Inside Convolutional Networks: Visualising Image Classification Models and Saliency Maps},
%   author={Simonyan, Karen and Vedaldi, Andrea and Zisserman, Andrew},
%   booktitle={International Conference on Learning Representations (ICLR) Workshop},
%   year={2014}
% }
%
% @inproceedings{shrikumar2017learning,
%   title={Learning Important Features Through Propagating Activation Differences},
%   author={Shrikumar, Avanti and Greenside, Peyton and Kundaje, Anshul},
%   booktitle={Proceedings of the 34th International Conference on Machine Learning (ICML)},
%   pages={3145--3153},
%   year={2017}
% }
%
% @inproceedings{springenberg2015striving,
%   title={Striving for Simplicity: The All Convolutional Net},
%   author={Springenberg, Jost Tobias and Dosovitskiy, Alexey and Brox, Thomas and Riedmiller, Martin},
%   booktitle={International Conference on Learning Representations (ICLR) Workshop},
%   year={2015}
% }
%
% @inproceedings{smilkov2017smoothgrad,
%   title={SmoothGrad: Removing Noise by Adding Noise},
%   author={Smilkov, Daniel and Thorat, Nikhil and Kim, Been and Vi{\'e}gas, Fernanda and Wattenberg, Martin},
%   booktitle={ICML Workshop on Visualization for Deep Learning},
%   year={2017}
% }
%
% @article{erion2021improving,
%   title={Improving Performance of Deep Learning Models with Axiomatic Attribution Priors and Expected Gradients},
%   author={Erion, Gabriel and Janizek, Joseph D. and Sturmfels, Pascal and Lundberg, Scott M. and Lee, Su-In},
%   journal={Nature Machine Intelligence},
%   volume={3},
%   number={7},
%   pages={620--631},
%   year={2021},
%   publisher={Nature Publishing Group}
% }
%
% @inproceedings{zeiler2014visualizing,
%   title={Visualizing and Understanding Convolutional Networks},
%   author={Zeiler, Matthew D. and Fergus, Rob},
%   booktitle={European Conference on Computer Vision (ECCV)},
%   pages={818--833},
%   year={2014},
%   organization={Springer}
% }
%
% @inproceedings{fong2017interpretable,
%   title={Interpretable Explanations of Black Boxes by Meaningful Perturbation},
%   author={Fong, Ruth C. and Vedaldi, Andrea},
%   booktitle={Proceedings of the IEEE International Conference on Computer Vision (ICCV)},
%   pages={3429--3437},
%   year={2017}
% }
%
% @inproceedings{fong2019understanding,
%   title={Understanding Deep Networks via Extremal Perturbations and Smooth Masks},
%   author={Fong, Ruth and Patrick, Mandela and Vedaldi, Andrea},
%   booktitle={Proceedings of the IEEE/CVF International Conference on Computer Vision (ICCV)},
%   pages={2950--2958},
%   year={2019}
% }
%
% @inproceedings{petsiuk2018rise,
%   title={RISE: Randomized Input Sampling for Explanation of Black-box Models},
%   author={Petsiuk, Vitali and Das, Abir and Saenko, Kate},
%   booktitle={British Machine Vision Conference (BMVC)},
%   year={2018}
% }
%
% @inproceedings{ribeiro2016why,
%   title={``Why Should I Trust You?'': Explaining the Predictions of Any Classifier},
%   author={Ribeiro, Marco Tulio and Singh, Sameer and Guestrin, Carlos},
%   booktitle={Proceedings of the 22nd ACM SIGKDD International Conference on Knowledge Discovery and Data Mining},
%   pages={1135--1144},
%   year={2016}
% }
%
% @inproceedings{lundberg2017unified,
%   title={A Unified Approach to Interpreting Model Predictions},
%   author={Lundberg, Scott M. and Lee, Su-In},
%   booktitle={Advances in Neural Information Processing Systems (NeurIPS)},
%   volume={30},
%   pages={4765--4774},
%   year={2017}
% }
%
% @article{strumbelj2014explaining,
%   title={Explaining Prediction Models and Individual Predictions with Feature Contributions},
%   author={{\v{S}}trumbelj, Erik and Kononenko, Igor},
%   journal={Knowledge and Information Systems},
%   volume={41},
%   number={3},
%   pages={647--665},
%   year={2014},
%   publisher={Springer}
% }
%
% @article{olson2021counterfactual,
%   title={Counterfactual State Explanations for Reinforcement Learning Agents via Generative Deep Learning},
%   author={Olson, Matthew L. and Khanna, Roli and Neal, Lawrence and Li, Fuxin and Wong, Weng-Keen},
%   journal={Artificial Intelligence},
%   volume={295},
%   pages={103455},
%   year={2021},
%   publisher={Elsevier}
% }
%
% @inproceedings{adebayo2018sanity,
%   title={Sanity Checks for Saliency Maps},
%   author={Adebayo, Julius and Gilmer, Justin and Muelly, Michael and Goodfellow, Ian and Hardt, Moritz and Kim, Been},
%   booktitle={Advances in Neural Information Processing Systems (NeurIPS)},
%   volume={31},
%   pages={9505--9515},
%   year={2018}
% }
%
% @inproceedings{abnar2020quantifying,
%   title={Quantifying Attention Flow in Transformers},
%   author={Abnar, Samira and Zuidema, Willem},
%   booktitle={Proceedings of the 58th Annual Meeting of the Association for Computational Linguistics (ACL)},
%   pages={4190--4197},
%   year={2020}
% }
%
% @inproceedings{bastani2018verifiable,
%   title={Verifiable Reinforcement Learning via Policy Extraction},
%   author={Bastani, Osbert and Pu, Yewen and Solar-Lezama, Armando},
%   booktitle={Advances in Neural Information Processing Systems (NeurIPS)},
%   volume={31},
%   pages={2494--2504},
%   year={2018}
% }
